\input{header}

% Pakiety do ladnego skladu pseudokodu (ladowane po headerze, wciaz w preambule).
% Jesli header.sty juz je laduje, ponowne \usepackage jest bezpieczne.
\usepackage{listings}
\usepackage{xcolor}

\lstdefinestyle{pseudo}{
    basicstyle=\ttfamily\small,
    keywordstyle=\bfseries,
    numbers=left,
    numberstyle=\tiny\color{gray},
    numbersep=8pt,
    frame=single,
    framesep=6pt,
    breaklines=true,
    captionpos=b,
    keywords={funkcja, jezeli, wpp, zwroc, dla, kazdego, dopoki, oraz, lub, nie}
}

% Bez wcięć akapitowych - akapity oddzielone odstępem.
\setlength{\parindent}{0pt}
\setlength{\parskip}{6pt plus 2pt minus 1pt}

\begin{document}
\frontpage{%
    Projektowanie i analiza algorytmów
}{%
    Informatyczne Systemy Automatyki
}{%
    środa TP $17^{05}-18^{45}$
}{-}{-}{%
    Jakub Włodarski 284184%
}{%
    https://github.com/kuba-wlo/Projektowanie-i-analiza-algorytmow-2%
}{15.06.2026 r}{%
    2
}
\pagestyle{empty}
\hypersetup{hidelinks}
\tableofcontents
\newpage
\pagestyle{plain}


\section{Wprowadzenie}

Kółko i krzyżyk należy do najprostszych gier o sumie zerowej z pełną informacją: obaj gracze
w każdym momencie znają cały stan rozgrywki, a wygrana jednego z nich oznacza przegraną
drugiego. W klasycznej odmianie partię toczy się na planszy $3\times3$, gdzie celem jest
ustawienie trzech własnych znaków w jednej linii. Przestrzeń stanów jest tu na tyle mała,
że grę da się rozwiązać w całości -- przy optymalnej grze obu stron partia zawsze kończy się
remisem.

Celem niniejszego projektu nie było jednak odtworzenie samego wariantu $3\times3$, lecz
zbudowanie jego \textbf{uogólnienia}, w którym to użytkownik decyduje o rozmiarze kwadratowej
planszy oraz o liczbie znaków w rzędzie wymaganej do zwycięstwa. Dzięki temu klasyczne
$3\times3$ staje się tylko jednym z wielu możliwych ustawień (program obsługuje plansze aż do
$20\times20$). Wraz ze wzrostem planszy przestrzeń stanów rośnie wykładniczo i gra przestaje
być rozwiązywalna metodą siłową -- to właśnie ten fakt nadaje projektowi charakter
algorytmiczny.

Przeciwnik komputerowy został oparty na algorytmie \textbf{MinMax z odcinaniem alfa-beta},
wzbogaconym o heurystyczną funkcję oceny, ograniczenie zbioru rozważanych pól oraz adaptacyjne
sterowanie głębokością przeszukiwania. Te trzy techniki wspólnie utrzymują grywalność również
na dużych planszach, gdzie samo przeszukiwanie pełnego drzewa byłoby nieobliczalne. Zgodnie
z listą tematów wybrano zadanie \emph{Kółko i krzyżyk} (ocena maksymalna 4,5), a dodatkowo
przygotowano wersję graficzną opartą na bibliotece Qt~6 (podnoszącą maksymalną ocenę do 5,0).

Całość napisano w języku \texttt{C++17}. Logikę gry oraz sztuczną inteligencję wydzielono
do osobnej biblioteki statycznej, z której korzystają dwie niezależne nakładki: graficzna
(Qt~6) oraz konsolowa (przeznaczona do szybkiego testowania).


\section{Opis gry}

\subsection{Uogólnione zasady}

Plansza jest kwadratem o boku $N$, gdzie $N \in \{3,\dots,20\}$. Każde pole może być puste albo
zajęte znakiem jednego z graczy ($\times$ lub $\circ$). Gracze wykonują ruchy naprzemiennie,
przy czym partię zawsze rozpoczyna znak $\times$. Zwycięża ten, kto jako pierwszy ustawi $k$
swoich znaków w nieprzerwanej linii, gdzie długość wygrywającej linii $k$ jest parametrem
konfigurowanym przez użytkownika i spełnia $3 \le k \le N$. Linia może biec w jednym z czterech
kierunków: poziomo, pionowo oraz po obu przekątnych. Jeśli wszystkie pola zostaną zajęte, a żaden
z graczy nie zbuduje wygrywającej linii, partia kończy się remisem.

Konfiguracja jest sanityzowana jeszcze przed rozpoczęciem rozgrywki: długość linii $k$ jest
zaciskana do przedziału $[1, N]$ (\texttt{std::clamp}), dzięki czemu niemożliwe jest zażądanie
linii dłuższej niż bok planszy, a samo GUI ogranicza wybór do sensownego zakresu $[3, N]$.

\subsection{Tryby gry i konfiguracja}

Program udostępnia następujące ustawienia rozgrywki:
\begin{itemize}
    \item rozmiar planszy $N$ (od 3 do 20),
    \item długość wygrywającej linii $k$ (od 3 do $N$),
    \item znak gracza: $\times$ (rozpoczyna) lub $\circ$ (wówczas partię rozpoczyna komputer),
    \item tryb gry: dwóch graczy przy jednym komputerze albo gra przeciwko sztucznej inteligencji,
    \item poziom trudności AI: łatwy, średni lub trudny,
    \item rozpoczęcie nowej partii z bieżącą konfiguracją.
\end{itemize}

Ponieważ pierwszy ruch zawsze należy do znaku $\times$, wybór znaku $\circ$ przez człowieka jest
równoznaczny z oddaniem otwarcia komputerowi.


\section{Architektura programu}

Aplikację rozbito na trzy moduły, co pozwala w pełni oddzielić logikę gry od warstwy prezentacji.
Wspólny rdzeń (logika + AI) zebrano w bibliotece statycznej \texttt{ttt\_core}; korzystają z niej
zarówno interfejs graficzny, jak i wersja konsolowa, dzięki czemu obie nakładki dzielą dokładnie
ten sam, raz przetestowany kod reguł i algorytmu.

\begin{table}[H]
\centering
\caption{Podział projektu na moduły}
\begin{tabular}{@{}lll@{}}
\toprule
\textbf{Katalog} & \textbf{Artefakt} & \textbf{Zawartość}\\
\midrule
\texttt{src/core/}    & biblioteka \texttt{ttt\_core} & logika gry oraz algorytm AI (bez GUI)\\
\texttt{src/gui/}     & \texttt{tictactoe\_gui}       & interfejs graficzny w Qt~6\\
\texttt{src/console/} & \texttt{tictactoe\_console}   & wersja tekstowa do testowania\\
\bottomrule
\end{tabular}
\end{table}

Najważniejsze klasy rdzenia i ich odpowiedzialności:
\begin{itemize}
    \item \texttt{Board} -- przechowuje stan kwadratowej planszy $N\times N$; zna wyłącznie
          rozmieszczenie znaków i nie wie nic o warunku zwycięstwa.
    \item \texttt{GameRules} -- wykrywa wygraną (zliczanie znaków w linii w czterech kierunkach)
          i wyznacza status gry; celowo oddzielona od planszy, by warunek zwycięstwa można było
          zmieniać niezależnie od reprezentacji pól.
    \item \texttt{Game} -- kontroler rozgrywki spinający planszę, zasady i kolejność ruchów;
          operuje wyłącznie na znakach, nie wie nic o graczach ani o GUI.
    \item \texttt{Player} -- abstrakcyjny interfejs gracza; jego polimorficznymi implementacjami
          są \texttt{HumanPlayer} (ruch dostarcza interfejs) oraz \texttt{AIPlayer}
          (ruch wylicza algorytm).
\end{itemize}

Takie rozwarstwienie sprawia, że nakładka jedynie odpytuje obiekt \texttt{Game} o stan i przekazuje
mu ruchy, natomiast cała logika decyzyjna pozostaje zamknięta w bibliotece. Wymiana interfejsu
(np. konsola zamiast Qt) nie wymaga żadnych zmian w regułach gry ani w AI.


\section{Logika gry i wykrywanie wygranej}

Sercem warstwy reguł jest klasa \texttt{GameRules}. Po każdym ruchu trzeba ustalić, czy nie
powstała właśnie wygrywająca linia. Zamiast przeglądać całą planszę, wykrywanie wygranej zostało
zrealizowane \textbf{inkrementalnie}: sprawdzane jest tylko otoczenie ostatnio postawionego znaku.
Dla każdego z czterech kierunków funkcja \texttt{countInLine} zlicza nieprzerwany ciąg znaków
przechodzący przez ostatnie pole -- w obie strony, włącznie z samym polem. Jeśli w którymkolwiek
kierunku długość ciągu osiągnie $k$, ruch kończy partię zwycięstwem.

Takie podejście jest istotnie tańsze niż pełny przegląd planszy: skoro tylko ostatni ruch mógł
utworzyć nową linię, wystarczy zbadać linie przechodzące właśnie przez to pole. Koszt pojedynczego
sprawdzenia jest rzędu $O(k)$ na kierunek, czyli $O(k)$ łącznie, zamiast $O(N^{2}k)$ przy naiwnym
skanowaniu całej planszy. Pełny przegląd (\texttt{hasWon}) pozostaje dostępny jako wariant
pomocniczy, używany m.in. wtedy, gdy nie jest znany ostatni ruch.


\section{Zastosowane techniki sztucznej inteligencji}

Logikę przeciwnika komputerowego zamknięto w klasie \texttt{AIPlayer}, która łączy kilka technik:
przeszukiwanie drzewa gry algorytmem MinMax z odcinaniem alfa-beta, ograniczenie zbioru
rozważanych pól, adaptacyjny dobór głębokości, heurystyczną ocenę pozycji oraz losowanie spośród
równorzędnych ruchów.

\subsection{Algorytm MinMax z odcinaniem alfa-beta}

MinMax to klasyczny algorytm wyznaczania optymalnego ruchu w grach o sumie zerowej. Przeszukuje
drzewo możliwych przebiegów partii, naprzemiennie maksymalizując wynik gracza AI i minimalizując
wynik przeciwnika (przy założeniu, że przeciwnik również gra optymalnie). Odcinanie alfa-beta
utrzymuje dwie granice -- $\alpha$ (najlepszy gwarantowany wynik gracza maksymalizującego) oraz
$\beta$ (najlepszy gwarantowany wynik minimalizującego) -- i przerywa analizę gałęzi, gdy
$\beta \le \alpha$, ponieważ taka gałąź nie może już wpłynąć na ostateczny wybór ruchu. Pozwala to
pominąć znaczną część drzewa bez zmiany wyniku.

Aby AI preferowało szybkie wygrane, od wartości wygranej odejmowana jest bieżąca głębokość:
natychmiastowa wygrana jest premiowana wyżej niż taka sama wygrana osiągnięta głębiej w drzewie.
Symetrycznie -- przegraną AI stara się odwlec jak najdłużej.

\begin{lstlisting}[style=pseudo, caption={Szkielet MinMaxa z odcinaniem alfa-beta (uproszczony).}, label={lst:minimax}]
funkcja minimax(plansza, glebokosc, alpha, beta, doRuchu, limit):
    jezeli plansza pelna: zwroc 0            # remis
    jezeli glebokosc >= limit: zwroc evaluate(plansza)
    maks = (doRuchu == AI)
    najlepszy = maks ? -INF : +INF
    dla kazdego ruchu w kandydaci(plansza):
        ustaw ruch
        jezeli ruch wygrywa:
            wynik = maks ? (WIN - glebokosc) : -(WIN - glebokosc)
        wpp:
            wynik = minimax(plansza, glebokosc+1, alpha, beta,
                            przeciwnik(doRuchu), limit)
        cofnij ruch
        jezeli maks: najlepszy = max(najlepszy, wynik); alpha = max(alpha, najlepszy)
        wpp:         najlepszy = min(najlepszy, wynik); beta  = min(beta,  najlepszy)
        jezeli beta <= alpha: zwroc najlepszy   # odciecie alfa-beta
    zwroc najlepszy
\end{lstlisting}

\subsection{Ograniczenie zbioru kandydatów}

Współczynnik rozgałęzienia drzewa gry równa się liczbie wolnych pól -- dla planszy $20\times20$
daje to na starcie aż $400$ możliwych ruchów, przez co pełne drzewo jest nieobliczalne. Aby temu
zapobiec, AI rozważa wyłącznie pola \emph{sąsiadujące} (w promieniu jednego pola) z już
postawionymi znakami. Ruchy w odległych, pustych obszarach nie mają znaczenia taktycznego, więc
ich pominięcie nie pogarsza jakości gry, a drastycznie zmniejsza rozgałęzienie. Na pustej planszy
zbiorem kandydatów jest samo pole środkowe. W końcówkach (do dziewięciu wolnych pól) rozważane są
wszystkie wolne pola, co umożliwia pełne, dokładne przeszukanie.

\subsection{Adaptacyjna głębokość przeszukiwania}

Nawet po ograniczeniu kandydatów drzewo na dużych planszach pozostaje kosztowne, dlatego głębokość
przeszukiwania dobierana jest do bieżącego rozgałęzienia: im więcej ruchów do rozważenia, tym płycej
schodzi przeszukiwanie. Do dziewięciu wolnych pól (małe plansze, końcówki) drzewo przeszukiwane jest
w całości, z dokładnością do limitu wynikającego z poziomu trudności. Po osiągnięciu limitu pozycja
jest oceniana funkcją heurystyczną, zamiast schodzić przeszukiwaniem do liści.

\begin{table}[H]
\centering
\caption{Górny limit głębokości w zależności od liczby kandydatów (rozgałęzienia)}
\begin{tabular}{@{}lc@{}}
\toprule
\textbf{Liczba kandydatów} & \textbf{Limit głębokości}\\
\midrule
$\le 6$  & 6\\
$\le 12$ & 4\\
$\le 24$ & 3\\
$> 24$   & 2\\
\bottomrule
\end{tabular}
\end{table}

\subsection{Heurystyczna funkcja oceny}

Gdy limit głębokości zostanie osiągnięty, pozycja pozostaje nierozstrzygnięta i trzeba ją oszacować.
Funkcja oceny przesuwa ,,okno'' długości $k$ po wszystkich liniach planszy (we wszystkich czterech
kierunkach) i klasyfikuje każde okno:
\begin{itemize}
    \item okno \emph{mieszane} (zawierające znaki obu graczy) jest martwe -- żaden gracz nie może
          go już dokończyć -- więc wnosi $0$ punktów,
    \item okno zawierające wyłącznie znaki AI wnosi punkty dodatnie,
    \item okno zawierające wyłącznie znaki przeciwnika wnosi punkty ujemne.
\end{itemize}
Wartość okna rośnie wykładniczo wraz z liczbą znaków (mnożnik $10$ na każdy kolejny znak), co mocno
premiuje linie bliskie ukończenia. Wartość jest ograniczona z góry (do $100\,000$), aby ocena
heurystyczna nigdy nie została pomylona z rzeczywistą wygraną (oznaczaną wartością $10^{6}$).

\subsection{Wybór otwarcia i poziomy trudności}

Poziom trudności steruje górnym limitem głębokości przeszukiwania (\texttt{maxDepth}). Na poziomie
łatwym AI patrzy w praktyce tylko na pozycję bezpośrednio po własnym ruchu (nie symuluje odpowiedzi
przeciwnika), natomiast na poziomie trudnym przeszukuje na tyle głęboko, że na planszy $3\times3$
gra idealnie.

\begin{table}[H]
\centering
\caption{Odwzorowanie poziomu trudności na limit głębokości}
\begin{tabular}{@{}lc@{}}
\toprule
\textbf{Poziom} & \textbf{\texttt{maxDepth}}\\
\midrule
łatwy   & 1\\
średni  & 4\\
trudny  & 9\\
\bottomrule
\end{tabular}
\end{table}

Od poziomu zależy także \emph{otwarcie}, gdy AI rusza jako pierwsze. Na poziomie trudnym AI zawsze
gra w środek planszy -- pole leżące na największej liczbie linii, a więc najlepsze otwarcie. Na
poziomie średnim losuje środek lub jedno z jego sąsiednich pól, a na łatwym -- dowolne pole. Dzięki
temu na niższych poziomach kolejne partie nie zaczynają się zawsze tak samo.

\subsection{Losowanie spośród równorzędnych ruchów}

Gdy kilka ruchów uzyskuje identyczną, najlepszą ocenę, AI losuje jeden z nich (generatorem
\texttt{std::mt19937}). Dzięki temu kolejne partie różnią się od siebie, mimo że gra pozostaje
optymalna. Aby pula ,,równorzędnych'' ruchów była rzetelna, każdy ruch w korzeniu drzewa oceniany
jest pełnym oknem $(-\infty, +\infty)$, bez zacieśniania $\alpha$ między kolejnymi ruchami.
W przeciwnym razie odcięcia alfa-beta zwracałyby dla części ruchów jedynie oszacowanie (granicę)
zamiast dokładnej wartości i do puli trafiałyby ruchy faktycznie gorsze. Odcięcia pozostają
natomiast aktywne w głębszych węzłach drzewa, gdzie służą wyłącznie przyspieszeniu.

\subsection{Analiza złożoności}

Złożoność czasowa czystego algorytmu MinMax wynosi $O(b^{d})$, gdzie $b$ to współczynnik
rozgałęzienia, a $d$ -- głębokość przeszukiwania. W przypadku optymistycznym (dobre uporządkowanie
ruchów) odcinanie alfa-beta redukuje ją do $O(b^{d/2})$, co odpowiada możliwości przeszukania
dwukrotnie głębszego drzewa w tym samym czasie. W projekcie obie składowe są dodatkowo kontrolowane:
ograniczenie kandydatów do sąsiedztwa zmniejsza efektywne $b$, a adaptacyjny dobór głębokości
zmniejsza $d$ tam, gdzie $b$ jest duże. Koszt pojedynczej oceny heurystycznej jest rzędu
$O(N^{2}k)$ -- dla każdego z $O(N^{2})$ pól startowych i czterech kierunków przeglądane jest okno
długości $k$. Z kolei inkrementalne wykrywanie wygranej po ruchu kosztuje jedynie $O(k)$, co czyni
je tanim również w węzłach przeszukiwania.


\section{Interfejs graficzny i wersja konsolowa}

\subsection{Wersja graficzna (Qt~6)}

Interfejs graficzny zbudowano w oparciu o moduł \texttt{Widgets} biblioteki Qt~6. Okno główne
udostępnia panel konfiguracji (rozmiar planszy, długość linii, znak gracza, poziom trudności,
tryb gry z AI), planszę oraz pasek statusu informujący o przebiegu partii. Plansza jest responsywna
-- skaluje się wraz z oknem, pozostając kwadratowa i wyśrodkowana. Kliknięcie w wolne pole wykonuje
ruch gracza, po czym w trybie z AI ruch wylicza \texttt{AIPlayer}. Widget planszy
(\texttt{BoardWidget}) jedynie odpytuje stan z obiektu \texttt{Game} i go rysuje, a o kliknięciach
informuje sygnałem \texttt{cellClicked}.

\begin{figure}[H]
\centering
\includegraphics[width=0.8\textwidth]{plansza.png}
\caption{Okno główne aplikacji -- panel konfiguracji i plansza.}
\label{fig:gui-konfig}
\end{figure}

\begin{figure}[H]
\centering
\includegraphics[width=0.8\textwidth]{gra.png}
\caption{Rozgrywka przeciwko sztucznej inteligencji.}
\label{fig:gui-gra}
\end{figure}

\subsection{Wersja konsolowa}

Wersja tekstowa korzysta z tej samej biblioteki \texttt{ttt\_core}, co interfejs graficzny.
Pozwala szybko sprawdzić zasady i zachowanie AI bez uruchamiania okna Qt: rysuje planszę z indeksami
wierszy i kolumn oraz wczytuje ruchy z klawiatury, walidując zakres wprowadzanych wartości. Była ona
podstawowym narzędziem do weryfikacji poprawności rdzenia jeszcze przed podłączeniem warstwy
graficznej.

\begin{figure}[H]
\centering
\includegraphics[width=0.60\textwidth]{konsola.png}
\caption{Wersja konsolowa -- plansza z indeksami pól.}
\label{fig:konsola}
\end{figure}


\section{Podsumowanie i wnioski}

W ramach projektu zaimplementowano uogólnioną grę w kółko i krzyżyk wraz z przeciwnikiem sterowanym
przez sztuczną inteligencję. Gracz może swobodnie definiować rozmiar planszy oraz długość
wygrywającej linii, dzięki czemu klasyczne $3\times3$ jest tylko jednym z możliwych wariantów.
Przeciwnik AI oparto na algorytmie MinMax z odcinaniem alfa-beta.

Najważniejszym wnioskiem płynącym z realizacji jest to, że sam MinMax -- choć gwarantuje optymalną
grę -- jest niepraktyczny dla większych plansz ze względu na wykładniczy wzrost drzewa gry.
Kluczowe okazało się połączenie go z trzema technikami ograniczającymi koszt obliczeń: zawężeniem
zbioru rozważanych pól do sąsiedztwa już postawionych znaków, adaptacyjnym doborem głębokości do
bieżącego rozgałęzienia oraz heurystyczną oceną pozycji w węzłach, do których przeszukiwanie nie
dochodzi. Dzięki nim AI pozostaje grywalne nawet na planszach $20\times20$, jednocześnie grając
idealnie na planszy $3\times3$ na najwyższym poziomie trudności.

Na większych planszach faktyczna głębokość przeszukiwania bywa ograniczana przez adaptację, przez co
różnica między poziomami trudności jest tam mniej wyraźna -- jest to świadomy kompromis między siłą
gry a czasem odpowiedzi. Wprowadzenie poziomów trudności oraz losowania spośród równorzędnych ruchów
istotnie poprawiło odbiór gry: partie się nie powtarzają, a przeciwnika można dopasować do
umiejętności gracza.

Możliwe kierunki dalszego rozwoju to wprowadzenie tablicy transpozycji (spamiętywanie ocenionych
pozycji), porządkowanie ruchów zwiększające skuteczność odcięć alfa-beta oraz pogłębianie iteracyjne
z limitem czasowym zamiast stałego limitu głębokości.


\begin{thebibliography}{9}
\bibitem{cormen} T.~H.~Cormen, C.~E.~Leiserson, R.~L.~Rivest, C.~Stein,
\emph{Wprowadzenie do algorytmów}, WNT, Warszawa, 2012.
\bibitem{russell} S.~Russell, P.~Norvig,
\emph{Artificial Intelligence: A Modern Approach}, 4th ed., Pearson, 2021
(rozdział o przeszukiwaniu drzew gry: MinMax i odcinanie alfa-beta).
\bibitem{knuth} D.~E.~Knuth, R.~W.~Moore,
\emph{An Analysis of Alpha-Beta Pruning}, Artificial Intelligence, vol.~6, no.~4, s.~293--326, 1975.
\bibitem{qt} Dokumentacja Qt~6 -- moduł Widgets, \url{https://doc.qt.io/}.
\bibitem{cppref} Dokumentacja biblioteki standardowej C++, \url{https://en.cppreference.com/}.
\end{thebibliography}

\end{document}
